\section{Latest Resource-Matched Coverage Analysis}
\label{app:latest-evidence}
This section and the main resource-matched tables use the frozen
\texttt{final\_closure\_v9} evidence. Earlier sections preserve the original
experiments, figures, and their support. They are not additional independent
replications of this panel, and their historical use of ``primary'' denotes
their own frozen analysis plan rather than replacing the present main summary.

\subsection{Source and coverage details}
% evidence: configs/final_robustness_v9.json; paper/generated/final_closure_v9/verification/verification.json
EmoEEG-MC uses snapshot 1.0.7; Urban Appraisal uses snapshot 1.0.0
\citep{openneuro_ds005540,openneuro_ds006850}.
Emo has 1,260 eligible trials from 30 participants and 42 stimulus identities.
Urban uses 3,471 valence trials from 62 participants and 56 images; other
scales are not copied into valence/arousal labels. A source recording with
confirmed truncation was excluded before prediction inspection, as documented
in the frozen coverage manifest. Urban's dataset-description license says CC0,
while its README says CC BY 4.0. We document both declarations and do not
redistribute the source recording or image material.

% evidence: src/openaffect_eeg/final_robustness.py; configs/final_robustness_v9.json
SHA-256 ordering with design seed 20260908 assigns participant and stimulus
identities to five blocks. Block $b$ supplies test pairs; block $b+1$ modulo five
supplies validation pairs; population training initially uses identities from
the other blocks. Test pairs lie on diagonal blocks only. Coverage means every
eligible participant and stimulus is tested once, not every possible trial.
The same test support is retained throughout the dose grid. Stimulus-label
injection may use participants from the validation block; only the tested
participants are categorically excluded from population fitting at all doses.
This is not a claim of fully disjoint validation identities after injection.

\subsection{Bounded model settings and preprocessing}
% evidence: configs/final_robustness_v9.json; scripts/prepare_urban_tensors_v9.py; src/openaffect_eeg/final_robustness.py
Ridge R1 selects alpha from $\{0.1,1,10,100\}$ and R2 from
$\{1,10,100,1000\}$. Both band-power and frozen LaBraM use the two grids.
The historical Emo feature key \texttt{de} selects channel log-bandpower; it
does not make this panel a separate differential-entropy implementation.
Standard EEGNet is Braindecode 1.2.0 with F1=8, D=2, F2=16, dropout 0.25,
AdamW weight decay $10^{-4}$, batch size 128, at most 100 epochs and patience
15. L1 and L2 use learning rates $10^{-3}$ and $10^{-4}$. Validation MAE
selects training length; the selected length is refitted on train plus
validation. Both direct and residual branches are retained. Calibration changes
post-fit offsets only; it is not a sweep of few-shot fine-tuning algorithms.

% evidence: configs/final_robustness_v9.json; scripts/prepare_urban_tensors_v9.py
Emo uses the existing 100-Hz raw trial archive and ten-second EEGNet crops.
Urban uses 64 EEG channels in physical volts, excludes ECG and GSR, applies
common-average reference and polyphase resampling from 500 to 100 Hz, and
uses the three-second image interval before the rating response. At evaluation,
window predictions are averaged within a trial. Downstream learned
normalization is fitted without test or calibration trials. Frozen LaBraM
uses the project's recorded checkpoint and channel/unit adaptation; its upstream
training-data identity provenance is not certified here.

\subsection{Complete primary and secondary summaries}
% evidence: paper/generated/evidence_argument_v10/verification.json
All tables below are generated from the same hashed aggregates as the main
tables. No new fitting, bootstrap, or model selection was performed to produce
this manuscript revision. The score identity $EC-P=(PC-P)+(EC-PC)$ is checked
for all 300 task--setting--dose combinations after block aggregation, with
maximum numerical error below $10^{-12}$. This validates a numerical identity,
not a causal decomposition. Paired intervals are copied from paired contrast
outputs, never constructed by subtracting separate confidence bounds.
\input{../generated/evidence_argument_v10/full_sensitivity.tex}

\subsection{Resource effects and statistical limits}
% evidence: paper/generated/final_closure_v9/normal_t/table_added_value_summary.csv
For the EEG-free prior, participant-dose changes from $(0,0)$ to $(8,0)$ are
$+0.2036$ $[+0.0955,+0.3116]$ on Emo and $+0.1564$ $[+0.0492,+0.2637]$ on
Urban. Stimulus-dose changes to $(0,8)$ are $+0.4907$ $[+0.3581,+0.6233]$ and
$+0.3408$ $[+0.2161,+0.4655]$. Joint changes to $(8,8)$ are $+0.6327$
$[+0.4810,+0.7844]$ and $+0.4181$ $[+0.2928,+0.5435]$. These are heuristic
t-normal sensitivity intervals; the corresponding percentile intervals also
exclude zero. They are not the same contrasts as $PC-P$ at a fixed stimulus
dose in the main endpoint table.

The bootstrap reuses crossed participant and stimulus multiplicities across
all fitted predictors. Averages are computed within block and then across
blocks rather than pretending repeated grid predictions are new observations.
The finite-cluster coverage diagnostics and invalid-draw counts constrain
interpretation. These are fixed-fit uncertainty summaries, not repeated-training
confidence statements or equivalence tests. Neither isolated pointwise
intervals nor dose-family bands imply simultaneous coverage over every method,
task, and secondary endpoint inspected during benchmark development.

\subsection{Known-truth evaluator stress tests}
\label{app:evaluator-validation}
% evidence: paper/generated/submission_closure_v11/evaluation_method_validation.json
We tested the estimator on outcomes with known construction before interpreting
the EEG results. In the v6 oracle suite, 16 settings with 100 replicates each
produced the following operating characteristics. ``Detection'' denotes a
strictly positive interval and is not a universal power guarantee.

\begin{center}
\small\setlength{\tabcolsep}{4pt}
\begin{tabularx}{\linewidth}{>{\raggedright\arraybackslash}p{0.19\linewidth}>{\raggedright\arraybackslash}p{0.23\linewidth}>{\raggedright\arraybackslash}X}
\toprule
Generator & Recorded behavior & Interpretation boundary \\
\midrule
Labels only & mean difference exactly zero & algebraic negative control only \\
Independent noise & 0/16 cells detected & finite evidence, not a proof of zero false positives \\
Trial signal & detection 0.85--1.00; coverage 0.92--1.00 & detects the injected trial signal in these settings \\
Participant constant & minimum coverage 0.82 & exposes finite-cluster weakness; no biological or causal interpretation \\
\bottomrule
\end{tabularx}
\end{center}
% evidence: paper/generated/submission_closure_v11/evaluation_method_validation.json

The v7 finite-cluster study planned 9,600 fits; 9,584 were valid and 16 were
rejected by the declared support rule. In the target $32\times24$ cluster
setting, percentile, basic, and t-normal coverages were 0.933, 0.930, and
0.957. Their minima in the deliberately small $8\times6$ setting were 0.837,
0.820, and 0.940. The often wider t-normal interval is therefore retained as
a sensitivity analysis, not promoted as a generally calibrated solution.
% evidence: paper/generated/submission_closure_v11/evaluation_method_validation.json

\subsection{Rule basis and three-state decisions}
% source: configs/audit_rule_catalog.json
The checker implements observable evidence rules rather than inferring
scientific validity from a score. Every rule returns allow, block, or
unverifiable; the last state is preserved when required provenance is absent.

\begin{center}
\small\setlength{\tabcolsep}{4pt}
\begin{tabularx}{\linewidth}{>{\raggedright\arraybackslash}p{0.18\linewidth}>{\raggedright\arraybackslash}p{0.27\linewidth}>{\raggedright\arraybackslash}X}
\toprule
Rule family & Required evidence & Example boundary \\
\midrule
Trial lineage & trial UIDs in every partition & repeated derived windows block trial-isolation claims \\
Identity holdout & participant/stimulus UIDs & overlap blocks only the identity claim it violates \\
Preprocess/selection & fit and selection UIDs & missing provenance is unverifiable, not a pass \\
Matched resources & predictor-specific label access & unequal calibration blocks an EEG-added-value claim \\
Fixed support & test-set hashes by comparator & unequal test support blocks a paired comparison \\
\bottomrule
\end{tabularx}
\end{center}
These rules operationalize documented leakage and benchmark-reporting concerns
\citep{brookshire2024leakage,lei2025leakage,bendrich2022feature,reuel2024betterbench};
they do not certify data truth, upstream pretraining history, or causal validity.

\section{Source-Linked Execution Cases and Reuse}
\label{app:execution-cases}
% evidence: paper/generated/final_closure_v9/external_audit/public_summary.json
The executed EEGain commit is
\texttt{d6892f5586181f345969651a9baedabd828bb1c5}; its
\texttt{EEGDataloader.split\_train\_val} is in
\texttt{eegain/data/loader.py}. LibEER is pinned to
\texttt{dddff9776dbdae21195fe320dff0a5ba61628a18}; we execute
\texttt{get\_split\_index}, \texttt{index\_to\_data}, and the empty-validation
branch at lines 100--102 of \texttt{LibEER/EEGNet\_train.py}.
File hashes, selected source ranges, case origin, and outputs are recorded in
\texttt{external\_audit/public\_summary.json} and its paired case table.
Marker windows are generated from real trial identities, not independent EEG
measurements. No external classifier score is inferred from the branch tests.

The same EEGain execution record produces different answers to different
contracts: its held-out test participants remain excluded, while an additional
requirement that validation trials be disjoint from training can be violated.
Changing the claim to what was actually isolated is one valid repair. When a
strict validation-trial contract is required, partitioning complete trials
before generating windows is the appropriate construction instead. LibEER's
explicit train/validation/test path passes the exercised split checks; its
empty-validation fallback should not be treated as independent validation.
The machine-readable repair record
\texttt{submission\_closure\_v11/external\_claim\_repair.json} verifies this
strict-block/scoped-allow pair across all five exercised EEGain folds while
holding the pinned function, trial universe, and execution output fixed. The
repair changes the supportable claim; it does not alter or re-estimate a model
score.

\begin{small}\begin{verbatim}
python scripts/run_workflow_contract_toy.py toy_contract
python -m openaffect_eeg.workflow_audit \
  --trials toy_contract/trials.tsv \
  --execution toy_contract/unmatched_calibration_execution.json \
  --contract toy_contract/unmatched_calibration_contract.json \
  --output toy_contract/checked.json
\end{verbatim}\end{small}
The mismatch example exits with the documented blocked status; the matched
example passes observable checks. Omitting preprocessing evidence preserves
uncertainty. This is a developer-operated executable example, not independent
human acceptance. The checker validates consistency of supplied execution
records, not adversarially falsified provenance.

\paragraph{AI-assisted development.}
AI assistance was used for protocol and code drafting, debugging, literature
triage, analysis support, and manuscript development. The recorded EEG
predictions, bootstrap summaries, and contract decisions are generated by
the documented numerical models and deterministic checks, not by an LLM
acting as an emotion predictor or audit judge. Automated checks and agent-run
examples are not independent human review. Authors remain responsible for
source verification, scientific interpretation, and submission approval.

\section{Historical Overview Figures}
% constants-source: paper/figure_plan_zh.json
The following approved figures retain earlier evidence without changing its
source version. The initial data layer covers the original four datasets;
Urban is the later source described above. These figures are not the latest
matched-increment tables and must not be read as latent-emotion decompositions.
\placeholderfigure{figures/figure1.pdf}{2in}{Original trial-first data and
evaluation contract. Stimulus identity is unavailable for ds006866. The later
Urban source and resource-matched predictor pairs are specified in the main
methods and latest-evidence appendix.}{fig:benchmark-contract}
\placeholderfigure{figures/figure2.pdf}{2.5in}{Original EmoEEG-MC and MusicEEG
total calibrated prediction surfaces on fixed trials. This is total system
CCC, not EEG added value. The new coverage analysis and matched contrasts are
reported separately in the main tables; MusicEEG retains its original feasible
grid rather than being relabelled as a full maximum-dose grid.}{fig:identity-exposure}
\placeholderfigure{figures/figure3.pdf}{2.3in}{Original representation probes.
Observed and permuted identity decoding describe available information,
not proof that an affect predictor uses identity or that identity is a pure
nuisance.}{fig:identity-audit}
